% Apuntes de teoría en PDF · 1.º ESO · Unidad 6: Álgebra y ecuaciones (versión en castellano)
% Mateix estil que els apunts de Batxillerat (franja i caixes de teoria en blau, exemples en roig), però més visuals,
% curts i amb dibuixos, per a l'alumnat de 3r ESO. Es compila amb pdflatex dues vegades (per als enllaços de l'índex),
% amb logo-fentmates.pdf (_eines/logo/pdf/horitzontal-clar.pdf) al costat.
\documentclass[11pt,a4paper]{article}
\usepackage[utf8]{inputenc}
\usepackage[T1]{fontenc}
\usepackage[spanish,es-noshorthands,es-nodecimaldot,es-nolists]{babel}
\usepackage[sfdefault,book,lining]{FiraSans}
\usepackage[italic,noplusnominus,noequal]{mathastext}
\linespread{1.18}
\usepackage[a4paper,top=2.4cm,bottom=2.2cm,left=1.6cm,right=1.6cm,headheight=1.1cm,headsep=0.5cm,footskip=1cm]{geometry}
\usepackage{amsmath,amssymb,cancel}
\usepackage{icomma}
\usepackage{graphicx}
\usepackage[table]{xcolor}
\usepackage{tikz}\usetikzlibrary{calc,positioning,arrows.meta,decorations.pathreplacing,shapes.geometric}
\usepackage{fontawesome5}
\usepackage{ccicons}
\usepackage{fancyhdr}
\usepackage{enumitem}
\usepackage{array}
\usepackage{needspace}
\usepackage[most]{tcolorbox}
\usepackage{titlesec}
\usepackage[hidelinks]{hyperref}
\hypersetup{pdftitle={Apuntes de teoría. Unidad 6: Álgebra y ecuaciones. 1.º ESO},pdfsubject={Matemáticas 1.º ESO},pdfkeywords={lenguaje algebraico, expresiones algebraicas, valor numérico, monomios, operaciones con monomios, ecuaciones, solución, ecuaciones equivalentes, transposición, paréntesis, denominadores, problemas}}

% colors de la web (bloc d'Àlgebra: blau)
\definecolor{fmnavy}{HTML}{1B2748}
\definecolor{fmblue}{HTML}{3D6AA0}
\definecolor{fmbluel}{HTML}{4E81BD}
\definecolor{fmyellow}{HTML}{FFC94D}
\definecolor{fmgold}{HTML}{E9A919}
\definecolor{fmaccent}{HTML}{FF7A50}
\definecolor{fmdef}{HTML}{1D4F73}\definecolor{fmdefbg}{HTML}{EEF4F9}
\definecolor{fmex}{HTML}{A3203A}\definecolor{fmexbg}{HTML}{FBF0F2}
\definecolor{fmsf}{HTML}{C2410C}
\definecolor{fmok}{HTML}{15803D}
\definecolor{fmmuted}{HTML}{5D6685}
\definecolor{fmsurface}{HTML}{EEF1F8}
\definecolor{fmA}{HTML}{2F6FD6}\definecolor{fmB}{HTML}{EA6A10}\definecolor{fmR}{HTML}{E11D48}\definecolor{fmK}{HTML}{16A34A}\definecolor{fmV}{HTML}{7C3AED}

% lletres
\newcommand{\comic}{\fontfamily{ComicNeue-TLF}\selectfont}
\newcommand{\heavy}{\fontseries{eb}\selectfont}
\newcommand{\semi}{\fontseries{sb}\selectfont}
\DeclareUnicodeCharacter{2212}{-}

\setlength{\parindent}{0pt}
\setlength{\parskip}{7pt}
\setlength{\lineskip}{4pt}\setlength{\lineskiplimit}{4pt}
\setlength{\emergencystretch}{3em}
\setlist{itemsep=3pt,topsep=3pt,leftmargin=1.4em}
\setcounter{secnumdepth}{1}
\renewcommand{\thesection}{6.\arabic{section}}
\titleformat{\section}[hang]{\color{fmnavy}\fontsize{23}{27}\heavy\selectfont}{\tcbox[on line,colback=fmblue,colframe=fmblue,arc=3mm,boxsep=2pt,left=6pt,right=6pt,top=3pt,bottom=3pt]{\color{white}\thesection}}{0.45em}{}[\vspace{5pt}{\color{fmblue}\titlerule[2pt]}]
\titlespacing*{\section}{0pt}{4pt}{6pt}
\titleformat{\subsection}{\color{fmblue}\Large\semi}{}{0pt}{\faCaretRight\ }
\titlespacing*{\subsection}{0pt}{8pt}{2pt}

% frase de cada apartat, «a mà»
\newcommand{\frase}[1]{\par\vspace{-2pt}{\comic\Large\color{fmbluel}\faCommentDots[regular]\ \ #1\par}\vspace{2pt}}
\newcommand{\seccio}[3]{\Needspace{16\baselineskip}\section{#2}\hypertarget{#1}{}\frase{#3}}
% notes a mà dins dels exemples
\newcommand{\nota}[1]{{\comic\small\color{fmex!85!black}#1}}
\newcommand{\notab}[1]{{\comic\small\color{fmblue}#1}}
% bé / malament
\newcommand{\be}[1]{{\color{fmok}\faCheckCircle}\ \ensuremath{#1}}
\newcommand{\mal}[1]{{\color{fmR}\faTimesCircle}\ {\color{fmR}\ensuremath{#1}}}
% pas numerat
\newcommand{\pas}[1]{\tikz[baseline=(p.base)]\node[circle,fill=fmblue,text=white,inner sep=1.2pt,minimum size=1.25em,font=\small\heavy](p){#1};}
% etiqueta de color
\newcommand{\chip}[2]{\tcbox[on line,colback=#1!14,colframe=#1!60,arc=2mm,boxsep=1pt,left=5pt,right=5pt,top=1pt,bottom=1pt]{#2}}
\newcommand{\fmkey}[1]{\colorbox{fmyellow!55}{\strut #1}}
\newcommand{\cA}[1]{{\color{fmA}#1}}\newcommand{\cB}[1]{{\color{fmB}#1}}\newcommand{\cK}[1]{{\color{fmK}#1}}\newcommand{\cR}[1]{{\color{fmR}#1}}\newcommand{\cV}[1]{{\color{fmV}#1}}

% caixes
\tcbset{fmbase/.style={enhanced,arc=3mm,boxrule=1.2pt,left=4mm,right=4mm,top=2mm,bottom=2.5mm,before skip=8pt,after skip=8pt,fonttitle=\heavy\large,before upper=\raggedright}}
\newtcolorbox{idea}[1][Idea clave]{fmbase,colback=fmdefbg,colframe=fmdef,coltitle=white,title={\faKey\ \ #1}}
\newtcolorbox{exemple}[1][Ejemplo]{fmbase,colback=fmexbg,colframe=fmex,colbacktitle=fmexbg,coltitle=fmex,title={\faPencil*\ \ #1},titlerule=0.5pt,titlerule style={fmex!40,dashed}}
\newtcolorbox{compte}[1][¡Cuidado!]{fmbase,colback=fmsf!6,colframe=fmsf,coltitle=white,fonttitle=\comic\bfseries\Large,title={\faExclamationTriangle\ \ #1}}
\newtcolorbox{truc}[1][Truco]{fmbase,colback=fmyellow!20,colframe=fmgold,colbacktitle=fmyellow,coltitle=fmnavy,fonttitle=\comic\bfseries\Large,title={\faLightbulb[regular]\ \ #1}}

% dibuix de la targeta de la unitat, en blanc
\newcommand{\fmcard}{\begin{tikzpicture}[x=0.062cm,y=-0.062cm,white,line width=3pt,line cap=round]
  \begin{scope}[transparency group,opacity=0.15]
  \node[anchor=base west,inner sep=0pt,font=\fontsize{30}{30}\selectfont] at (14,44) {$2x+3=11$};
  \node[anchor=base west,inner sep=0pt,font=\fontsize{26}{26}\selectfont] at (14,100) {$x=4$};
  \end{scope}
\end{tikzpicture}}
\newcommand{\fmbrand}{\raisebox{-0.25\height}{\includegraphics[height=0.5cm]{logo-fentmates.pdf}}}

% regla de Ruffini: #1 = a, #2 = coeficients, #3 = fila del mig (sense el primer), #4 = fila de baix, #5 = nombre de coeficients
\newcommand{\ruffini}[5]{\begin{tikzpicture}[x=1.25cm,y=0.85cm,every node/.style={font=\large}]
  \foreach \v [count=\i] in {#2} \node (t\i) at (\i,0) {$\v$};
  \foreach \v [count=\i] in {#3} {\pgfmathtruncatemacro\j{\i+1}\node (m\j) at (\j,-1) {$\v$};}
  \foreach \v [count=\i] in {#4} \node (b\i) at (\i,-2.1) {$\v$};
  \node[text=fmB] at (0,-1) {$#1$};
  \draw[line width=1.1pt,fmnavy] (0.45,0.5)--(0.45,-1.5) (0.45,-1.5)--(#5+0.5,-1.5);
  \node[draw=fmR,line width=1.3pt,rounded corners=2mm,minimum size=0.75cm] at (#5,-2.1) {};
  \draw[-{Latex[length=2.2mm]},fmK,line width=1.1pt] (t1.south) -- (b1.north);
  \pgfmathtruncatemacro\last{#5-1}
  \foreach \i in {1,...,\last} {\pgfmathtruncatemacro\j{\i+1}
    \draw[-{Latex[length=2.2mm]},fmB,line width=1.1pt,shorten >=1pt] (b\i.north east) -- (m\j.south west);
    \draw[-{Latex[length=2.2mm]},fmA!70,line width=1pt,dashed,shorten >=2pt] ([xshift=8pt]t\j.south) -- ([xshift=8pt]b\j.north);}
\end{tikzpicture}}

% capçalera i peu
\pagestyle{fancy}\fancyhf{}
\fancyhead[L]{\small\fmbrand}
\fancyhead[R]{\small\color{fmmuted}1.º ESO · Unidad 6: Álgebra y ecuaciones}
\fancyfoot[L]{\scriptsize\color{fmmuted}\href{https://creativecommons.org/licenses/by-nc-sa/4.0/deed.es}{\ccbyncsa}\ \ FentMates · fentmates.org · CC BY-NC-SA 4.0}
\fancyfoot[R]{\small\color{fmmuted}\thepage}
\renewcommand{\headrulewidth}{0.4pt}\renewcommand{\headrule}{\color{fmmuted!40}\hrule height \headrulewidth}
\fancypagestyle{primera}{\fancyhf{}\renewcommand{\headrulewidth}{0pt}
  \fancyfoot[L]{\scriptsize\color{fmmuted}\href{https://creativecommons.org/licenses/by-nc-sa/4.0/deed.es}{\ccbyncsa}\ \ FentMates · fentmates.org · CC BY-NC-SA 4.0}
  \fancyfoot[R]{\small\color{fmmuted}\thepage}}

% targeta de l'índex
\newcommand{\itile}[4]{\begin{tcolorbox}[enhanced,colback=white,colframe=fmblue!35,arc=3mm,boxrule=1pt,boxsep=0pt,left=6pt,right=6pt,top=5pt,bottom=5pt,before skip=0pt,after skip=7pt]
  \hyperlink{#1}{\begin{minipage}[c]{0.9cm}\centering{\color{fmblue}\LARGE #3}\end{minipage}\hspace{6pt}%
  \begin{minipage}[c]{0.78\linewidth}{\color{fmblue}\heavy #2}\\[-2pt]{\color{fmnavy}\semi\large #4}\end{minipage}}\end{tcolorbox}}

% taula amb vores arrodonides: el contingut es retalla amb la forma (sense @{} a les vores ni entre columnes)
\newsavebox\tabclipbox
\newcommand{\taulaclip}[3][3mm]{\sbox\tabclipbox{#3}\begin{tikzpicture}[baseline=(tcn.base)]
  \begin{scope}\clip[rounded corners=#1] (0,-\dp\tabclipbox) rectangle (\wd\tabclipbox,\ht\tabclipbox);
  \node[inner sep=0pt,outer sep=0pt,anchor=base west] (tcn) at (0,0) {\usebox\tabclipbox};\end{scope}
  \draw[#2,rounded corners=#1] (0,-\dp\tabclipbox) rectangle (\wd\tabclipbox,\ht\tabclipbox);
\end{tikzpicture}}
\newcommand{\taularod}[3][fmblue]{\taulaclip{#1,line width=1.2pt}{\renewcommand{\arraystretch}{1.5}\begin{tabular}{#2}#3\end{tabular}}}
% balança
\newcommand{\balanca}[3]{\begin{tikzpicture}[scale=#3,transform shape,every node/.style={anchor=center}]
  \fill[fmnavy] (-0.45,-1.9) rectangle (0.45,-1.75);
  \fill[fmnavy!80] (-0.12,-1.75) rectangle (0.12,0);
  \fill[fmgold] (0,0.05) circle (0.13);
  \draw[fmnavy,line width=2.4pt,line cap=round] (-2,0) -- (2,0);
  \foreach \s in {-1,1} {
    \draw[fmmuted,line width=0.8pt] (\s*2,0) -- (\s*2-0.75,-1.05) (\s*2,0) -- (\s*2+0.75,-1.05);
    \fill[fmyellow!70,draw=fmgold,line width=1pt] (\s*2-0.95,-1.05) arc[start angle=180,end angle=360,x radius=0.95,y radius=0.32] -- cycle;}
  \node[font=\Large,fill=fmA!15,draw=fmA,rounded corners=1.5mm,inner sep=3pt] at (-2,-0.68) {$#1$};
  \node[font=\Large,fill=fmB!15,draw=fmB,rounded corners=1.5mm,inner sep=3pt] at (2,-0.68) {$#2$};
\end{tikzpicture}}
% eixos amb quadrícula (dins d'un tikzpicture): xmin, xmax, ymin, ymax
\newcommand{\eixos}[4]{\draw[fmmuted!22,line width=0.4pt] (#1,#3) grid (#2,#4);
  \draw[-{Latex[length=2mm]},fmmuted,line width=0.8pt] (#1-0.2,0)--(#2+0.4,0) node[right,font=\small]{$x$};
  \draw[-{Latex[length=2mm]},fmmuted,line width=0.8pt] (0,#3-0.2)--(0,#4+0.4) node[above,font=\small]{$y$};}
% taula de traduccions (diccionari) amb les vores arrodonides
\newcommand{\dic}[2]{#1 & {\large$#2$}\\}
\newcommand{\dicratio}{0.62}
\newcommand{\dichead}{\begin{tabular}{@{}>{\raggedright}p{\dimexpr\dicratio\linewidth-12pt}>{\centering\arraybackslash}p{\dimexpr\linewidth-\dicratio\linewidth-12pt}@{}}\faComment[regular]\ Lo decimos así\dots & \dots\ así\end{tabular}}
\newenvironment{dicbox}{\begin{tcolorbox}[enhanced,arc=3mm,boxrule=1.2pt,colframe=fmblue,colback=white,coltitle=white,fonttitle=\heavy,
  title={\dichead},left=0pt,right=0pt,top=0pt,bottom=1.5mm,boxsep=0pt,lefttitle=6pt,righttitle=6pt,before skip=0pt,after skip=0pt]\rowcolors{1}{white}{fmsurface}%
  \begin{tabular}{@{\hspace{6pt}}>{\raggedright}p{\dimexpr\dicratio\linewidth-12pt}>{\centering\arraybackslash}p{\dimexpr\linewidth-\dicratio\linewidth-12pt}@{\hspace{6pt}}}}
  {\end{tabular}\end{tcolorbox}}
% resum
\newcommand{\rtile}[3]{\begin{tcolorbox}[enhanced,colback=white,colframe=fmblue!45,arc=3mm,boxrule=1pt,left=3mm,right=3mm,top=2.5mm,bottom=2.5mm,fontupper=\large,
  title={\tcbox[on line,colback=fmblue,colframe=fmblue,arc=2mm,boxsep=0pt,left=4pt,right=4pt,top=1pt,bottom=1pt]{\color{white}\small\heavy #1}\ \ #2},coltitle=fmnavy,colbacktitle=fmdefbg,fonttitle=\semi,toptitle=0.5mm,bottomtitle=0.5mm]#3\end{tcolorbox}}
% portada: #1 subtítol a mà
\newcommand{\portada}[1]{\thispagestyle{primera}
\begin{tikzpicture}[remember picture,overlay]
  \fill[fmblue] (current page.north west) rectangle ([yshift=-9.4cm]current page.north east);
  \begin{scope}
    \clip (current page.north west) rectangle ([yshift=-9.4cm]current page.north east);
    \node[anchor=north west,inner sep=0pt] at ([xshift=6.6cm,yshift=-1.3cm]current page.north west) {\fmcard};
  \end{scope}
  \node[anchor=north west,inner sep=0pt] at ([xshift=1.6cm,yshift=-1.1cm]current page.north west) {\tcbox[colback=white,colframe=white,arc=2.5mm,boxsep=0pt,left=7pt,right=7pt,top=5pt,bottom=5pt]{\includegraphics[height=0.75cm]{logo-fentmates.pdf}}};
\end{tikzpicture}
\vspace*{0.55cm}
\tcbox[colback=white!15!fmblue,colframe=white!15!fmblue,arc=2.5mm,boxsep=0pt,left=6pt,right=6pt,top=3pt,bottom=3pt]{\color{white}\small\bfseries Apuntes de teoría · 1.º ESO · Matemáticas}

\vspace{2pt}
{\raggedright\color{white}\fontsize{30}{36}\heavy\selectfont{\color{fmyellow}Unidad 6}\\[2pt]Álgebra y ecuaciones\par}
\vspace{8pt}
{\color{white}\comic\Large #1\par}
\vspace*{\dimexpr 7.4cm-\pagetotal\relax}
}
\newcommand{\llegenda}{\par\vspace{16pt}
{\color{fmnavy}\heavy\Large\faPalette\ \ Cómo se leen estos apuntes}\par\smallskip
\begin{tcbraster}[raster columns=4,raster equal height,raster column skip=8pt,raster before skip=12pt]
\tcbox[colback=fmdefbg,colframe=fmdef,arc=3mm,boxrule=1.2pt,halign=center,width=\linewidth]{{\color{fmdef}\Large\faKey}\\[2pt]{\color{fmdef}\heavy Idea clave}\\[-1pt]{\small lo esencial}}
\tcbox[colback=fmexbg,colframe=fmex,arc=3mm,boxrule=1.2pt,halign=center,width=\linewidth]{{\color{fmex}\Large\faPencil*}\\[2pt]{\color{fmex}\heavy Ejemplo}\\[-1pt]{\small paso a paso}}
\tcbox[colback=fmsf!6,colframe=fmsf,arc=3mm,boxrule=1.2pt,halign=center,width=\linewidth]{{\color{fmsf}\Large\faExclamationTriangle}\\[2pt]{\color{fmsf}\heavy ¡Cuidado!}\\[-1pt]{\small errores típicos}}
\tcbox[colback=fmyellow!20,colframe=fmgold,arc=3mm,boxrule=1.2pt,halign=center,width=\linewidth]{{\color{fmgold!80!black}\Large\faLightbulb[regular]}\\[2pt]{\color{fmnavy}\heavy Truco}\\[-1pt]{\small para recordarlo}}
\end{tcbraster}}
\newcommand{\resumtitol}{\clearpage
\hypertarget{resum}{}
{\color{fmnavy}\fontsize{23}{27}\heavy\selectfont\faClipboardList\ \ En resumen}\par\vspace{5pt}{\color{fmblue}\titlerule[2pt]}
\frase{Toda la unidad en una página.}}
% els 5 passos per a resoldre problemes: #2 = text del pas 2
\newcommand{\passosproblema}[2]{\begin{center}
\begin{tikzpicture}[s/.style={draw=fmblue,fill=fmdefbg,rounded corners=3mm,line width=1.2pt,minimum width=2.75cm,minimum height=2.45cm,align=center,font=\small},
  ic/.style={font=\LARGE,text=fmblue}]
  \foreach \i/\ic/\t in {1/\faBookOpen/{\textbf{Leemos}\\y entendemos},2/\faQuestionCircle/{#1},3/\faEquals/{#2},4/\faCogs/{\textbf{Resolvemos}},5/\faCheckDouble/{\textbf{Comprobamos}\\y respondemos}} {
    \node[s] (s\i) at (3.3*\i,0) {};
    \node[ic] at ([yshift=-0.5cm]s\i.north) {\ic};
    \node[align=center,font=\small,anchor=north] at ([yshift=-0.95cm]s\i.north) {\t};
    \node[circle,fill=fmblue,text=white,font=\small\heavy,inner sep=1.5pt,minimum size=1.3em] at (s\i.north west) {\i};}
  \foreach \i [evaluate=\i as \j using int(\i+1)] in {1,...,4} \draw[-{Latex[length=3mm]},fmmuted,line width=1.5pt] (s\i) -- (s\j);
\end{tikzpicture}
\end{center}}

\newcommand{\lo}{\mathopen{]}}\newcommand{\ro}{\mathclose{[}}
% punts ple i buit, nombres dels eixos i taula de valors
\newcommand{\pple}[2][fmblue]{\fill[#1] #2 circle (3.6pt);}
\newcommand{\nums}[4]{\foreach \k in {#1,...,#2} {\ifnum\k=0\else\node[below,font=\scriptsize,text=fmmuted] at (\k,0) {$\k$};\fi}
  \foreach \k in {#3,...,#4} {\ifnum\k=0\else\node[left,font=\scriptsize,text=fmmuted] at (0,\k) {$\k$};\fi}}
\newcommand{\tv}[3][fmblue]{\taulaclip{#1,line width=1.2pt}{\renewcommand{\arraystretch}{1.45}%
  \begin{tabular}{>{\columncolor{#1}\color{white}\heavy}c*{12}{w{c}{0.8cm}}}#2\\\rowcolor{#1!8}\cellcolor{#1}#3\end{tabular}}}
% etiqueta d'una recta o corba
\newcommand{\etq}[3]{\node[fill=white,inner sep=1.5pt,text=#1,font=\small] at #2 {#3};}


\newcommand{\passos}[1]{\begin{tcolorbox}[fmbase,colback=fmdefbg,colframe=fmdef,coltitle=white,title={\faListOl\ \ Cómo se resuelven}]#1\end{tcolorbox}}
\newcommand{\tsq}[2]{\filldraw[fill=fmA!30,draw=fmA,line width=0.8pt,rounded corners=0.5pt] (#1,#2) rectangle ++(0.9,0.9); \node[font=\small] at ($(#1,#2)+(0.45,0.45)$) {$x^2$};}
\newcommand{\minipar}[2]{\begin{tikzpicture}[x=0.55cm,y=0.42cm]
  \draw[fmmuted,-{Latex[length=2mm]},line width=0.8pt] (-2.3,0) -- (2.3,0) node[right,font=\scriptsize]{$X$};
  \draw[#2,line width=1.6pt,domain=-1.75:1.75,samples=40] plot (\x,{\x*\x+#1});
\end{tikzpicture}}
\def\intvdot#1#2#3{\ifx#2o\filldraw[fill=white,draw=#3,line width=1.3pt] (#1,0) circle (3.3pt);\fi\ifx#2c\fill[#3] (#1,0) circle (3.5pt);\fi}
\newcommand{\intv}[5]{\begin{tikzpicture}[x=0.4cm,baseline=-3pt]
  \draw[fmmuted,line width=0.8pt] (-6,0)--(7,0);
  \foreach \k in {-5,...,6} \draw[fmmuted,line width=0.6pt] (\k,-0.08)--(\k,0.08);
  \ifx#3a\def\xa{-6}\else\def\xa{#1}\fi
  \ifx#4a\def\xb{7}\else\def\xb{#2}\fi
  \ifx#4a\draw[#5,line width=2.6pt,-{Latex[length=2.6mm]}] (\xa,0)--(\xb,0);\else\ifx#3a\draw[#5,line width=2.6pt,{Latex[length=2.6mm]}-] (\xa,0)--(\xb,0);\else\draw[#5,line width=2.6pt] (\xa,0)--(\xb,0);\fi\fi
  \intvdot{#1}{#3}{#5}\intvdot{#2}{#4}{#5}
  \ifx#3a\else\node[below=2pt,font=\scriptsize] at (#1,0) {$#1$};\fi
  \ifx#4a\else\node[below=2pt,font=\scriptsize] at (#2,0) {$#2$};\fi
\end{tikzpicture}}


\begin{document}
\portada{Letras en lugar de números, y el número que hace cierta la igualdad.}

{\color{fmnavy}\heavy\Large\faMapSigns\ \ ¿Qué aprenderás?}\par\smallskip
\begin{minipage}[t]{0.485\linewidth}
\itile{sec1}{6.1}{\faFont}{Lenguaje algebraico}
\itile{sec2}{6.2}{\faEquals}{Expresiones. Valor numérico}
\itile{sec3}{6.3}{\faCubes}{Monomios}
\itile{sec4}{6.4}{\faPlus}{Operaciones con monomios}
\itile{sec5}{6.5}{\faBalanceScale}{Ecuaciones: elementos y solución}
\end{minipage}\hfill
\begin{minipage}[t]{0.485\linewidth}
\itile{sec6}{6.6}{\faIcon{exchange-alt}}{Resolución de ecuaciones}
\itile{sec7}{6.7}{\faIcon{layer-group}}{Ecuaciones con paréntesis}
\itile{sec8}{6.8}{\faDivide}{Ecuaciones con denominadores}
\itile{sec9}{6.9}{\faLightbulb}{Problemas con ecuaciones}
\itile{resum}{\phantom{6.9}}{\faClipboardList}{En resumen}
\end{minipage}
\llegenda

% ============================================================ 6.1
\clearpage
\seccio{sec1}{Lenguaje algebraico}{Usamos letras para los números que no conocemos.}

\vfill
\begin{minipage}[c]{0.27\linewidth}\centering
\begin{tikzpicture}
  \fill[fmyellow!75,draw=fmgold,line width=1.2pt,rounded corners=1.5mm] (0,0) rectangle (2.2,1.8);
  \fill[fmyellow!95!fmgold,draw=fmgold,line width=1.2pt,rounded corners=1mm] (-0.15,1.8) rectangle (2.35,2.25);
  \fill[fmaccent] (1.0,0) rectangle (1.2,2.25);
  \draw[fmaccent,line width=2pt] (1.1,2.25) .. controls (0.6,2.8) and (0.4,2.3) .. (1.1,2.25) .. controls (1.8,2.8) and (1.6,2.3) .. (1.1,2.25);
  \node[font=\fontsize{34}{34}\heavy\selectfont,text=fmnavy] at (0.52,0.9) {?};
  \node[font=\fontsize{34}{34}\heavy\selectfont,text=fmnavy] at (1.7,0.9) {?};
  \node[font=\fontsize{30}{30}\selectfont,text=fmblue] at (3.0,1.0) {$=x$};
\end{tikzpicture}
\end{minipage}\hfill
\begin{minipage}[c]{0.69\linewidth}
{\Large La letra $x$ es como una \textbf{caja} con un \textbf{número escondido}.}\par
Puede valer 2, 7, $-3$\dots{} Las letras que hacen de números se llaman \textbf{variables} (o \textbf{incógnitas}, si queremos averiguar cuánto valen).
\begin{idea}
Una \textbf{expresión algebraica} mezcla \cA{\textbf{números}}, \cB{\textbf{letras}} y \cK{\textbf{operaciones}}:\par\vspace{2pt}
{\centering\chip{fmblue}{$3x-5$}\ \ \chip{fmblue}{$2a+b^2$}\ \ \chip{fmblue}{$\dfrac{x+1}{2}$}\par}
\end{idea}
\end{minipage}

\vfill
\subsection{De palabras a letras}
\providecommand{\dic}{}\renewcommand{\dic}[2]{#1 & {\large$#2$}\\}
% taula del diccionari dins d'una caixa amb les vores arrodonides (la capçalera és el títol de la caixa)
\providecommand{\dichead}{}\renewcommand{\dichead}{\begin{tabular}{@{}>{\raggedright}p{\dimexpr0.66\linewidth-12pt}>{\centering\arraybackslash}p{\dimexpr0.34\linewidth-12pt}@{}}\faComment[regular]\ Lo decimos así\dots & \dots\ así\end{tabular}}
\makeatletter\@ifundefined{dicbox}{}{\let\dicbox\relax\let\enddicbox\relax}\makeatother\newenvironment{dicbox}{\begin{tcolorbox}[enhanced,arc=3mm,boxrule=1.2pt,colframe=fmblue,colback=white,coltitle=white,fonttitle=\heavy,
  title={\dichead},left=0pt,right=0pt,top=0pt,bottom=1.5mm,boxsep=0pt,lefttitle=6pt,righttitle=6pt,before skip=0pt,after skip=0pt]\rowcolors{1}{white}{fmsurface}%
  \begin{tabular}{@{\hspace{6pt}}>{\raggedright}p{\dimexpr0.66\linewidth-12pt}>{\centering\arraybackslash}p{\dimexpr0.34\linewidth-12pt}@{\hspace{6pt}}}}
  {\end{tabular}\end{tcolorbox}}
\begin{minipage}[t]{0.485\linewidth}
\begin{dicbox}
\dic{El doble de un número}{2x}
\dic{El triple de un número}{3x}
\dic{La mitad de un número}{\frac{x}{2}}
\dic{El cuadrado de un número}{x^2}
\dic{El cubo de un número}{x^3}
\dic{El opuesto de un número}{-x}
\end{dicbox}
\end{minipage}\hfill
\begin{minipage}[t]{0.485\linewidth}
\begin{dicbox}
\dic{El siguiente de un número}{x+1}
\dic{El anterior de un número}{x-1}
\dic{Un número par}{2x}
\dic{Un número impar}{2x+1}
\dic{La suma de dos números}{x+y}
\dic{El producto de dos números}{xy}
\end{dicbox}
\end{minipage}

\vfill
\begin{minipage}[t]{0.42\linewidth}
\begin{truc}
No escribimos el punto:\\ \mbox{$3\cdot x=3x$}. \ Ni el 1: \mbox{$1x=x$}.
\end{truc}
\end{minipage}\hfill
\begin{minipage}[t]{0.55\linewidth}
\begin{compte}[¡Cuidado con el orden!]
«El doble de un número, más 3» \ $\to$\ \ \be{2x+3}\\[2pt]
«El doble de (un número más 3)» \ $\to$\ \ \be{2(x+3)}
\end{compte}
\end{minipage}

\vfill
\subsection{Valor numérico}
{\large Cambiamos la letra por el número y hacemos las operaciones.} \ Es como una \textbf{máquina}: entra un número y sale otro.

\vspace{-10pt}
\begin{center}
\begin{tikzpicture}[scale=0.9,transform shape]
  \node[draw=fmblue,fill=fmdefbg,rounded corners=4mm,minimum width=4.6cm,minimum height=1.6cm,line width=1.4pt,font=\Large] (m) {$3x^2-5x+1$};
  \node[text=fmblue] at ([xshift=5mm,yshift=-4mm]m.north west) {\faCog};
  \node[text=fmblue] at ([xshift=-5mm,yshift=4mm]m.south east) {\faCog};
  \node[left=1.6cm of m,fill=fmyellow!60,rounded corners=2mm,font=\Large,inner sep=6pt] (in) {$x=-2$};
  \node[right=1.6cm of m,fill=fmK!20,draw=fmK,rounded corners=2mm,font=\Large\heavy,inner sep=6pt] (out) {$23$};
  \draw[-{Latex[length=3.5mm]},line width=1.8pt,fmmuted] (in) -- node[above,font=\comic\small]{entra} (m);
  \draw[-{Latex[length=3.5mm]},line width=1.8pt,fmmuted] (m) -- node[above,font=\comic\small]{sale} (out);
\end{tikzpicture}
\end{center}

\vfill
\begin{minipage}[t]{0.53\linewidth}
\begin{exemple}[Ejemplo: \ $3x^2-5x+1$ \ para \ $x=-2$]
{\large$3\cdot\cR{(-2)}^2-5\cdot\cR{(-2)}+1=12+10+1=\boxed{23}$}\\[2pt]
\nota{sustituimos (con paréntesis) y operamos}
\end{exemple}
\end{minipage}\hfill
\begin{minipage}[t]{0.44\linewidth}
\begin{compte}
¡Negativos, \textbf{siempre entre paréntesis}!\\[3pt]
\be{(-2)^2=4}\qquad \mal{-2^2=-4}\\[1pt]
{\small (son cosas distintas)}
\end{compte}
\end{minipage}

\vfill
% ============================================================ 6.2
\clearpage
\seccio{sec2}{Expresiones algebraicas. Valor numérico}{Sustituimos las letras por números y operamos.}

\vfill
\begin{idea}[Expresión algebraica]
Números y letras unidos por los signos de las operaciones: \ $2x+3$, \ $ab$, \ $x^2-5x+6$, \ $\dfrac{x+1}{2}$, \ $3(a-b)$.\\[3pt]
Cada sumando es un \textbf{término}: $x^2-5x+6$ tiene tres términos, $x^2$, $-5x$ y $6$. Las letras son las \textbf{variables}.
\end{idea}

\vfill
\begin{minipage}[c]{0.52\linewidth}
\begin{idea}[Valor numérico]
El número que obtenemos cuando \textbf{sustituimos} las letras por números y hacemos las operaciones, con la jerarquía de siempre.
\end{idea}
\end{minipage}\hfill
\begin{minipage}[c]{0.45\linewidth}
\begin{exemple}[Ejemplo: \ $2x+3$ \ para \ $x=4$]
$2\cdot4+3=8+3=\boxed{11}$
\end{exemple}
\end{minipage}

\vfill
\begin{compte}[Dos errores típicos]
Entre el número y la letra hay una \textbf{multiplicación}: si $x=4$, \ $3x=3\cdot4=12$, \ \mal{34}.\\[3pt]
Un negativo, siempre \textbf{entre paréntesis}: si $x=-3$, \ $x^2=(-3)^2=9$, \ \mal{-3^2=-9}.
\end{compte}

\vfill
\begin{exemple}[Ejemplos]
\begin{minipage}[t]{0.48\linewidth}
$x^2-5x+6$ \ para \ $x=2$:\\ $4-10+6=\boxed{0}$\\[4pt]
$3(a-b)$ \ para \ $a=7$, $b=2$:\\ $3\cdot5=\boxed{15}$\\[4pt]
$\dfrac{x+1}{2}$ \ para \ $x=9$: \ $\dfrac{10}{2}=\boxed{5}$
\end{minipage}\hfill
\begin{minipage}[t]{0.49\linewidth}
$5-2x$ \ para \ $x=-3$:\\ $5-2\cdot(-3)=5+6=\boxed{11}$\\[4pt]
$a^2+b^2$ \ para \ $a=-2$, $b=3$:\\ $4+9=\boxed{13}$\\[4pt]
$2l+2a$ \ para \ $l=8$, $a=5$: \ $\boxed{26}$ \ \nota{un perímetro}
\end{minipage}
\end{exemple}

\vfill
\begin{exemple}[Ejemplo: \ la fórmula del taxi, \ $P=3+1{,}2\cdot k$ \ ($k$: kilómetros)]
$k=5$: \ $P=3+6=\boxed{9\text{ €}}$ \qquad $k=12$: \ $P=3+14{,}4=\boxed{17{,}40\text{ €}}$ \qquad \nota{el valor numérico responde preguntas reales}
\end{exemple}

\vfill
\begin{truc}[Haz una tabla]
Para ver cómo cambia una expresión, calcula el valor numérico para varios valores:\\[3pt]
{\centering\small\setlength{\tabcolsep}{6pt}\begin{tabular}{c|ccccc}$x$&0&1&2&3&4\\\hline $2x+1$&1&3&5&7&9\end{tabular}\par}
\end{truc}

\vfill
% ============================================================ 6.3
\clearpage
\seccio{sec3}{Monomios}{Un monomio es un número multiplicado por letras.}

\vfill
\begin{center}
\begin{tikzpicture}
  \node[font=\fontsize{48}{48}\selectfont,inner sep=2pt] (c) {$\cA{-5}$};
  \node[font=\fontsize{48}{48}\selectfont,inner sep=2pt,right=2pt of c] (l) {$\cB{x^{\cK{2}}y^{\cK{1}}}$};
  \node[below left=0.9cm and -0.3cm of c,fill=fmA!14,draw=fmA,rounded corners=2mm,text=fmA,font=\heavy\large] (cl) {coeficiente};
  \node[below right=0.9cm and -0.6cm of l,fill=fmB!14,draw=fmB,rounded corners=2mm,text=fmB,font=\heavy\large] (ll) {parte literal};
  \node[right=2.4cm of l,fill=fmK!14,draw=fmK,rounded corners=2mm,text=fmK!80!black,font=\large,align=center] (g) {{\heavy grado}\\ $2+1=\mathbf{3}$};
  \draw[-{Latex[length=3mm]},fmA,line width=1.3pt] (cl) -- (c);
  \draw[-{Latex[length=3mm]},fmB,line width=1.3pt] (ll) -- (l);
  \draw[-{Latex[length=3mm]},fmK,line width=1.3pt] (g.west) to[bend right=25] ([yshift=-2pt]l.north east);
  \node[font=\comic,text=fmmuted,above=0pt of g] {suma los exponentes};
  \node[font=\comic\small,text=fmmuted,below=2pt of g] {si no hay exponente, es 1: \ $y=y^1$};
\end{tikzpicture}
\end{center}

\vfill
\begin{minipage}[c]{0.43\linewidth}
\begin{idea}
\begin{itemize}[leftmargin=1em]
\item \cA{\textbf{Coeficiente}}: el número (con el signo).
\item \cB{\textbf{Parte literal}}: las letras y los exponentes.
\item \cK{\textbf{Grado}}: la suma de los exponentes.
\end{itemize}
\end{idea}
\end{minipage}\hfill
\begin{minipage}[c]{0.54\linewidth}\centering\renewcommand{\arraystretch}{1.35}
\taularod{*{4}{>{\hspace{2pt}}c<{\hspace{2pt}}}}{
\rowcolor{fmblue}\color{white}\heavy Monomio & \color{white}\heavy Coeficiente & \color{white}\heavy Parte literal & \color{white}\heavy Grado\\
$7x^4$ & \cA{$7$} & \cB{$x^4$} & \cK{4}\\
\rowcolor{fmsurface}$-a^2b^3$ & \cA{$-1$} & \cB{$a^2b^3$} & \cK{$2+3=5$}\\
$x$ & \cA{$1$} & \cB{$x$} & \cK{1}\\
\rowcolor{fmsurface}$-6$ & \cA{$-6$} & \cB{---} & \cK{0}
}
\end{minipage}

\vfill
\begin{minipage}[t]{0.44\linewidth}
\begin{truc}
¿No ves el coeficiente? Es un \textbf{1}:\\
$x^2=\cA{1}x^2$ \qquad $-x=\cA{-1}x$\\[2pt]
Un número solo es un monomio de \textbf{grado 0}.
\end{truc}
\end{minipage}\hfill
\begin{minipage}[t]{0.53\linewidth}
\begin{compte}[Esto NO son monomios]
\mal{\dfrac{3}{x}}\ \ {\small la letra divide}\\[3pt]
\mal{\sqrt{x}}\ \ {\small hay una raíz}\\[3pt]
\mal{2x+3}\ \ {\small es una suma (un polinomio)}
\end{compte}
\end{minipage}

\vfill
\subsection{Monomios semejantes: la misma familia}
{\large Son \textbf{semejantes} si tienen \textbf{la misma parte literal}} (mismas letras y mismos exponentes). Es como tener el mismo apellido.

\vfill
\begin{center}
\begin{tikzpicture}[card/.style={draw=#1,fill=#1!10,line width=1.2pt,rounded corners=3mm,minimum width=2.6cm,minimum height=1.25cm,font=\Large}]
  \node[card=fmK] (a) at (0,0) {$3\cB{x^2y}$};
  \node[card=fmK] (b) at (3.4,0) {$-\frac12\cB{x^2y}$};
  \node[font=\Large,text=fmK] at (1.7,0) {\faUsers};
  \node[text=fmK!80!black,font=\comic\large] at (1.7,-1.05) {\faCheckCircle\ sí son semejantes};
  \node[card=fmR] (c) at (8.4,0) {$3\cB{x^2y}$};
  \node[card=fmR] (d) at (11.8,0) {$3\cB{xy^2}$};
  \node[font=\Large,text=fmR] at (10.1,0) {$\neq$};
  \node[text=fmR,font=\comic\large] at (10.1,-1.05) {\faTimesCircle\ no lo son (exponentes cambiados)};
\end{tikzpicture}
\end{center}
\begin{idea}[¿Para qué sirve?]
\centering\Large Solo se pueden \textbf{sumar y restar} los monomios \textbf{semejantes}.
\end{idea}

\vfill
% ============================================================ 6.4
\clearpage
\seccio{sec4}{Operaciones con monomios}{Para sumar, deben ser de la misma familia. Para multiplicar, no hace falta.}

\vfill
\subsection{Suma y resta}
\begin{minipage}[c]{0.48\linewidth}
\begin{tikzpicture}[font=\Large]
  \node at (0,0) {3 \textcolor{fmR}{\faIcon{apple-alt}}};
  \node at (1.15,0) {$+$};
  \node at (2.3,0) {2 \textcolor{fmR}{\faIcon{apple-alt}}};
  \node at (3.45,0) {$=$};
  \node at (4.6,0) {\textbf{5} \textcolor{fmR}{\faIcon{apple-alt}}};
  \node[text=fmok] at (5.7,0) {\faCheckCircle};
  \node at (0,-1.1) {3 \textcolor{fmR}{\faIcon{apple-alt}}};
  \node at (1.15,-1.1) {$+$};
  \node at (2.3,-1.1) {2 \textcolor{fmK}{\faLemon}};
  \node at (3.45,-1.1) {$\to$};
  \node[font=\comic\normalsize,text=fmR,align=center] at (4.75,-1.1) {no se pueden\\juntar};
  \node[text=fmR] at (5.95,-1.1) {\faTimesCircle};
\end{tikzpicture}
\end{minipage}\hfill
\begin{minipage}[c]{0.5\linewidth}
\begin{tikzpicture}[x=1cm,y=1cm]
  \foreach \i in {0,...,4} {\tsq{\i*1}{0}}
  \node[font=\Large] at (5.25,0.45) {$+$};
  \foreach \i in {0,...,2} {\tsq{5.6+\i*1}{0}}
  \node[font=\Large,anchor=west] at (0,-0.65) {$5x^2+3x^2=\mathbf{8x^2}$};
  \node[font=\comic,text=fmmuted,anchor=west] at (4.6,-0.65) {8 piezas iguales};
\end{tikzpicture}
\end{minipage}

\vfill
\begin{minipage}[t]{0.52\linewidth}
\begin{idea}
\textbf{Sumamos los coeficientes}. La parte literal se queda \textbf{igual}.\\[4pt]
$5x^2+3x^2=(5+3)x^2=8x^2$\\[2pt]
$4ab-7ab=-3ab$\\[2pt]
$5x^2+3x$ \ \notab{no son semejantes: se queda así}
\end{idea}
\end{minipage}\hfill
\begin{minipage}[t]{0.45\linewidth}
\begin{compte}[Error típico]
\mal{5x^2+3x^2=8x^4}\\[3pt]
\be{5x^2+3x^2=8x^2}\\[3pt]
Cuando \textbf{sumamos}, los exponentes \textbf{no cambian}.
\end{compte}
\end{minipage}

\vfill
\subsection{Producto, cociente y potencia}
\begin{tcbraster}[raster columns=3,raster equal height,raster column skip=8pt]
\begin{tcolorbox}[fmbase,colback=fmdefbg,colframe=fmdef,coltitle=white,title={\faTimes\ \ Producto},before upper=\centering]
Números: \textbf{multiplicamos}\\ Exponentes: \fmkey{\textbf{sumamos}}\\[3pt]
{\large$x^m\cdot x^n=x^{m+n}$}\\[4pt]
$(-3x^2y)\cdot(4xy^3)=-12x^3y^4$
\end{tcolorbox}
\begin{tcolorbox}[fmbase,colback=fmdefbg,colframe=fmdef,coltitle=white,title={\faDivide\ \ Cociente},before upper=\centering]
Números: \textbf{dividimos}\\ Exponentes: \fmkey{\textbf{restamos}}\\[3pt]
{\large$x^m:x^n=x^{m-n}$}\\[4pt]
$15x^5y^2:(3x^2y)=5x^3y$
\end{tcolorbox}
\begin{tcolorbox}[fmbase,colback=fmdefbg,colframe=fmdef,coltitle=white,title={\faSuperscript\ \ Potencia},before upper=\centering]
Número: \textbf{lo elevamos}\\ Exponentes: \fmkey{\textbf{multiplicamos}}\\[3pt]
{\large$(x^m)^n=x^{m\cdot n}$}\\[4pt]
$(2x^3)^2=4x^6$
\end{tcolorbox}
\end{tcbraster}

\vfill
\begin{minipage}[t]{0.62\linewidth}
\begin{exemple}
\vspace{-6pt}
\begin{align*}
2x^3\cdot5x^4&=(\cA{2\cdot5})\,x^{\cK{3+4}}=\boxed{10x^7}\\
20x^6y^3:(-4x^2y)&=(\cA{20:(-4)})\,x^{\cK{6-2}}y^{\cK{3-1}}=\boxed{-5x^4y^2}
\end{align*}
\nota{\cA{números con números} \quad\cK{exponentes de cada letra}}
\end{exemple}
\end{minipage}\hfill
\begin{minipage}[t]{0.35\linewidth}
\begin{truc}
Recuerda: \ $x^0=1$\\[4pt]
$6x^2y:(2x^2)=3\,\cR{x^0}\,y$\\[2pt]
$\phantom{6x^2y:(2x^2)}=3\cdot\cR{1}\cdot y=3y$
\end{truc}
\end{minipage}

\vfill
% ============================================================ 6.5
\clearpage
\seccio{sec5}{Ecuaciones. Elementos y solución}{Una ecuación es una pregunta: ¿qué número hace cierta la igualdad?}

\vfill
\begin{minipage}[c]{0.4\linewidth}\centering
\balanca{2x+3}{11}{0.95}
\end{minipage}\hfill
\begin{minipage}[c]{0.57\linewidth}
{\Large Una \textbf{ecuación} es como una\\ \textbf{balanza} en equilibrio.}\par
Los dos lados del signo $=$ tienen que «pesar» lo mismo. Hay que encontrar cuánto vale la $x$ para que pase.\\[2pt]
Aquí, \ $x=4$, \ porque \ $2\cdot4+3=11$.
\end{minipage}

\vspace{10pt}

\vfill
\begin{tcbraster}[raster columns=2,raster equal height,raster column skip=10pt]
\begin{tcolorbox}[fmbase,colback=fmK!8,colframe=fmK,coltitle=white,title={\faCheckCircle\ \ Identidad}]
Se cumple para \textbf{cualquier} valor de la letra.\par\smallskip
{\centering\chip{fmK}{$2(x+3)=2x+6$}\ \ \chip{fmK}{$(x+1)^2=x^2+2x+1$}\par}
\end{tcolorbox}
\begin{tcolorbox}[fmbase,colback=fmdefbg,colframe=fmdef,coltitle=white,title={\faSearch\ \ Ecuación}]
Solo se cumple para \textbf{algunos} valores (o para ninguno).\par\smallskip
{\centering\chip{fmblue}{$2x+3=11$}\ \ solo si \ $x=4$\par}
\end{tcolorbox}
\end{tcbraster}

\vfill
\subsection{Las partes de una ecuación}
\begin{center}
\begin{tikzpicture}[t/.style={font=\fontsize{28}{28}\selectfont,inner sep=2pt,text height=22pt,text depth=6pt},
  lb/.style={rounded corners=2mm,font=\large,align=center,anchor=south}]
  \node[t] (t1) {$3x^2$};
  \node[t,right=3pt of t1] (t2) {$-5$};
  \node[t,right=8pt of t2] (eq) {$=$};
  \node[t,right=8pt of eq] (t3) {$2x$};
  \node[t,right=3pt of t3] (t4) {$+3$};
  \draw[decorate,decoration={brace,amplitude=6pt,mirror},fmA,line width=1.2pt] ([yshift=-3pt]t1.south west) -- ([yshift=-3pt]t2.south east) node[midway,below=8pt,font=\large\heavy,text=fmA] {1.er miembro};
  \draw[decorate,decoration={brace,amplitude=6pt,mirror},fmB,line width=1.2pt] ([yshift=-3pt]t3.south west) -- ([yshift=-3pt]t4.south east) node[midway,below=8pt,font=\large\heavy,text=fmB] {2.º miembro};
  % punts on apunten les fletxes
  \coordinate (px) at ([xshift=23pt,yshift=-7pt]t1.north west);   % la x de 3x^2
  \coordinate (pe) at ([xshift=-5pt,yshift=5pt]t1.north east);    % l'exponent 2
  \coordinate (p5) at ([xshift=4pt,yshift=-5pt]t2.north);         % el 5
  \coordinate (p3) at ([xshift=5pt,yshift=-5pt]t4.north);         % el 3
  \node[lb,fill=fmV!12,draw=fmV] (inc) at ([xshift=-3.3cm,yshift=0.9cm]px) {{\heavy\color{fmV}incógnita}: $x$\\[-2pt]{\comic\small la letra}};
  \draw[-{Latex[length=3mm]},fmV,line width=1.3pt] (inc.east) to[out=0,in=150] (px);
  \node[lb,fill=fmK!12,draw=fmK] (gr) at ([xshift=-0.6cm,yshift=1.45cm]pe) {{\heavy\color{fmK!80!black}grado 2}\\[-2pt]{\comic\small el exponente mayor}};
  \draw[-{Latex[length=3mm]},fmK,line width=1.3pt] (gr.south -| pe) -- (pe);
  \node[lb,fill=fmB!12,draw=fmB] (ti) at ([xshift=4.4cm,yshift=1.45cm]p5) {{\heavy\color{fmB}términos independientes}\\[-2pt]{\comic\small los que no tienen letra}};
  \draw[-{Latex[length=3mm]},fmB,line width=1.3pt] (ti.south west) to[out=-150,in=80] (p5);
  \draw[-{Latex[length=3mm]},fmB,line width=1.3pt] (ti.south -| p3) -- (p3);
  \node[anchor=west,font=\comic,text=fmmuted,align=left] at ([xshift=1.1cm,yshift=-0.55cm]t4.east) {cada sumando\\es un \textbf{término}};
\end{tikzpicture}
\end{center}

\vfill
\begin{minipage}[t]{0.5\linewidth}
\begin{idea}[Solución]
Una \textbf{solución} es un valor de la $x$ que hace \textbf{cierta} la igualdad.\\[3pt]
\textbf{Resolver} una ecuación es encontrar \textbf{todas} las soluciones (puede tener una, varias o ninguna).
\end{idea}
\end{minipage}\hfill
\begin{minipage}[t]{0.47\linewidth}
\begin{compte}[¡Cuidado con el grado!]
Antes de mirar el grado, \textbf{operamos}:\\[3pt]
$x(x+2)=x^2+5$\\[2pt]
$\cancel{x^2}+2x=\cancel{x^2}+5\ \to\ 2x=5$\\[2pt]
Parecía de grado 2, pero es de \textbf{grado 1}.
\end{compte}
\end{minipage}

\vfill
\begin{exemple}[Ejemplo: \ ¿identidad o ecuación?]
\begin{tabular}{@{}p{0.47\linewidth}p{0.47\linewidth}@{}}
$3(x-2)=3x-6$\ \ $\to$\ \ \chip{fmK}{identidad} &
$3x-2=7$\ \ $\to$\ \ \chip{fmblue}{ecuación} \ \nota{solo si $x=3$}\\[6pt]
$(x-1)^2=x^2-2x+1$\ \ $\to$\ \ \chip{fmK}{identidad} &
$x^2=9$\ \ $\to$\ \ \chip{fmblue}{ecuación} \ \nota{solo si $x=3$ o $x=-3$}
\end{tabular}\\[4pt]
\nota{Para saberlo, operamos cada miembro: si quedan exactamente iguales, es una identidad.}
\end{exemple}

% ------------------------------------------------------------ 5.1 (solució i transposició)
\clearpage
\subsection{¿Es solución? Lo comprobamos}
{\large Cambiamos la $x$ por el número en \textbf{cada miembro por separado}.} Si los dos resultados coinciden, es solución.

\vfill
\begin{center}
\begin{tikzpicture}[card/.style={draw=#1,fill=#1!8,line width=1.2pt,rounded corners=3mm,minimum width=6.6cm,minimum height=2.4cm,align=center}]
  \node[font=\Large,fill=fmyellow!45,rounded corners=2mm,inner sep=5pt] at (4.1,1.85) {$x^2+1=3x-1$};
  \node[card=fmK] (a) at (0.6,0) {};
  \node[font=\large] at (0.6,0.6) {\chip{fmK}{$x=\mathbf{2}$}};
  \node at (0.6,-0.15) {$\cA{2^2+1=5}\qquad \cB{3\cdot2-1=5}$};
  \node[text=fmK!80!black,font=\comic\large] at (0.6,-0.8) {\faCheckCircle\ iguales: sí es solución};
  \node[card=fmR] (b) at (7.6,0) {};
  \node[font=\large] at (7.6,0.6) {\chip{fmR}{$x=\mathbf{3}$}};
  \node at (7.6,-0.15) {$\cA{3^2+1=10}\qquad \cB{3\cdot3-1=8}$};
  \node[text=fmR,font=\comic\large] at (7.6,-0.8) {\faTimesCircle\ distintos: no es solución};
\end{tikzpicture}
\end{center}

\vfill
\subsection{Ecuaciones equivalentes: la transposición}
{\large Para resolver, hacemos \textbf{lo mismo en los dos lados}: la balanza sigue en equilibrio.} Así obtenemos ecuaciones \textbf{equivalentes} (con las mismas soluciones) cada vez más sencillas, hasta que la $x$ queda \textbf{sola}.

\vfill
\begin{center}
\begin{tikzpicture}[t/.style={font=\fontsize{24}{24}\selectfont,inner sep=2pt,text height=19pt,text depth=6pt},
  l/.style={font=\comic,align=left,anchor=west}]
  % columna del signe =
  \node[t] (a3) at (0,0) {$=$};
  \node[t,anchor=east] (a2) at ([xshift=-6pt]a3.west) {$\cB{+5}$};
  \node[t,anchor=east] (a1) at ([xshift=-2pt]a2.west) {$3x$};
  \node[t,anchor=west] (a4) at ([xshift=6pt]a3.east) {$17$};
  \node[t] (b3) at (0,-1.7) {$=$};
  \node[t,anchor=east] (b1) at (a1.east |- b3) {$\cK{3}x$};
  \node[t,anchor=west] (b4) at ([xshift=6pt]b3.east) {$17\cB{-5}=12$};
  \node[t] (c3) at (0,-3.4) {$=$};
  \node[t,anchor=east] (c1) at (b1.east |- c3) {$x$};
  \node[t,anchor=west] (c4) at ([xshift=6pt]c3.east) {$12:\cK{3}=\boxed{4}$};
  \coordinate (m5) at ([xshift=48pt,yshift=-4pt]b4.north west);   % el -5
  \coordinate (d3) at ([xshift=34pt,yshift=-4pt]c4.north west);   % el 3 que divideix
  \draw[-{Latex[length=3mm]},fmB,line width=1.5pt] ([yshift=2pt]a2.south) to[out=-60,in=110] (m5);
  \draw[-{Latex[length=3mm]},fmK,line width=1.5pt] ([xshift=-1pt,yshift=2pt]b1.south) to[out=-60,in=110] (d3);
  \node[l,text=fmB] at (5.2,-0.85) {el $+5$ estaba \textbf{sumando}:\\pasa \textbf{restando}};
  \node[l,text=fmK!80!black] at (5.2,-2.55) {el $3$ estaba \textbf{multiplicando}:\\pasa \textbf{dividiendo}};
  \node[anchor=west] at (9.3,-1.5) {\balanca{3x+5}{17}{0.9}};
  \node[font=\comic,text=fmmuted,align=center] at (12.1,-3.05) {hacemos lo mismo\\en los dos platos};
\end{tikzpicture}
\end{center}

\vfill
\begin{tcbraster}[raster columns=2,raster equal height,raster column skip=10pt]
\begin{tcolorbox}[fmbase,colback=fmdefbg,colframe=fmdef,coltitle=white,title={\faPlus\,\faMinus\ \ Regla de la suma}]
Lo que está \textbf{sumando} en un miembro pasa al otro \textbf{restando}, y al revés.\par\smallskip
{\centering$x\cB{+7}=10\ \to\ x=10\cB{-7}=3$\par}
\end{tcolorbox}
\begin{tcolorbox}[fmbase,colback=fmdefbg,colframe=fmdef,coltitle=white,title={\faTimes\,\faDivide\ \ Regla del producto}]
Lo que está \textbf{multiplicando} a todo un miembro pasa al otro \textbf{dividiendo}, y al revés.\par\smallskip
{\centering$\cK{5}x=20\ \to\ x=\dfrac{20}{\cK{5}}=4$\par}
\end{tcolorbox}
\end{tcbraster}

\vfill
\begin{minipage}[t]{0.55\linewidth}
\begin{compte}[Cambia la operación, no el signo]
En \ $\cK{-2}x=8$, \ el $\cK{-2}$ pasa dividiendo \textbf{con su signo}:\\[4pt]
\be{x=\dfrac{8}{-2}=-4}\qquad \mal{x=\dfrac{8}{2}=4}
\end{compte}
\end{minipage}\hfill
\begin{minipage}[t]{0.42\linewidth}
\begin{truc}[Compruébalo siempre]
Sustituye la solución en la ecuación \textbf{del principio}:\\[4pt]
$3\cdot\cR{4}+5=17$ \ \be{}
\end{truc}
\end{minipage}

\vfill
% ============================================================ 6.6
\clearpage
\seccio{sec6}{Ecuaciones equivalentes. Resolución}{La $x$ solo aparece elevada a 1. Siempre tienen una solución.}

\vfill
\begin{idea}
{\centering\Large Las reducimos hasta \ \fmkey{$ax=b$} \ y despejamos: \quad $x=\dfrac{b}{a}$\par}
\end{idea}

\vfill
\begin{minipage}[c]{0.4\linewidth}\raggedright
\pas{1} Quitamos \textbf{paréntesis}.\\[6pt]
\pas{2} \textbf{Transponemos}: las $x$ a un lado y los números al otro.\\[6pt]
\pas{3} \textbf{Reducimos}: sumamos los semejantes hasta $ax=b$.\\[6pt]
\pas{4} \textbf{Despejamos} la $x$ y simplificamos.\\[6pt]
\pas{5} \textbf{Comprobamos}.
\end{minipage}\hfill
\begin{minipage}[c]{0.57\linewidth}
\begin{exemple}[Ejemplo: \ $3(x-2)-2(x+1)=4-x$]
\vspace{-6pt}
\begin{alignat*}{2}
&3x-6-2x-2=4-x &\qquad& \text{\nota{\pas{1}\ paréntesis}}\\
&3x-2x\cB{+x}=4\cB{+6+2} && \text{\nota{\pas{2}\ transponemos}}\\
&2x=12 && \text{\nota{\pas{3}\ reducimos}}\\
&x=\dfrac{12}{2}=\boxed{6} && \text{\nota{\pas{4}\ despejamos}}
\end{alignat*}
\nota{\pas{5}\ $3\cdot4-2\cdot7=-2$ \ y \ $4-6=-2$ \ \be{}}
\end{exemple}
\end{minipage}

\vfill
\begin{compte}[El menos delante de un paréntesis]
Cambia el signo de \textbf{todos} los términos de dentro:\\[4pt]
\be{-2(x+1)=-2x\,\cR{-}\,2}\qquad\quad \be{-(3-x)=-3\,\cR{+}\,x}\qquad\quad \mal{-2(x+1)=-2x+2}
\end{compte}

\vfill
\subsection{Casos especiales: cuando las $x$ desaparecen}
{\large A veces las $x$ se anulan y llegamos a \ $0x=b$.}

\vfill
\begin{tcbraster}[raster columns=2,raster equal height,raster column skip=10pt]
\begin{tcolorbox}[fmbase,colback=fmR!6,colframe=fmR,coltitle=white,title={\faTimesCircle\ \ $0x=$ un número: ninguna solución}]
$2(x+1)=2x+5$\\[2pt]
$2x+2=2x+5\ \to\ \cR{0x=3}$\\[4pt]
\notab{Ningún número multiplicado por 0 da 3.}
\end{tcolorbox}
\begin{tcolorbox}[fmbase,colback=fmK!8,colframe=fmK,coltitle=white,title={\faInfinity\ \ $0x=0$: infinitas soluciones}]
$3(x+2)-x=2x+6$\\[2pt]
$2x+6=2x+6\ \to\ \cK{0x=0}$\\[4pt]
\notab{Todos los números la cumplen: es una identidad.}
\end{tcolorbox}
\end{tcbraster}

\vfill
\begin{minipage}[t]{0.5\linewidth}
\begin{exemple}[Ejemplo: \ $7x-4=3x+8$]
\vspace{-6pt}
\begin{alignat*}{2}
&7x-3x=8+4 &\qquad& \text{\nota{transponemos}}\\
&4x=12\ \to\ x=\dfrac{12}{4}=\boxed{3}
\end{alignat*}
\nota{Comprobación: $7\cdot3-4=17$ \ y \ $3\cdot3+8=17$ \be{}}
\end{exemple}
\end{minipage}\hfill
\begin{minipage}[t]{0.47\linewidth}
\begin{exemple}[Ejemplo: \ $2(x+5)=3x-1$]
\vspace{-6pt}
\begin{alignat*}{2}
&2x+10=3x-1\\
&2x-3x=-1-10\\
&-x=-11\ \to\ \boxed{x=11} &\quad& \text{\nota{cambiamos el signo}}
\end{alignat*}
\end{exemple}
\end{minipage}

\vfill
% ============================================================ 6.7
\clearpage
\seccio{sec7}{Ecuaciones con paréntesis}{Primero quitamos los paréntesis; después, como siempre.}

\vfill
\begin{minipage}[c]{0.45\linewidth}
\passos{\pas{1} Quitamos \textbf{paréntesis} (propiedad distributiva).\\[3pt]
\pas{2} \textbf{Transponemos}: las $x$ a un lado, los números al otro.\\[3pt]
\pas{3} \textbf{Reducimos}.\\[3pt]
\pas{4} \textbf{Despejamos} la $x$.\\[3pt]
\pas{5} \textbf{Comprobamos}.}
\end{minipage}\hfill
\begin{minipage}[c]{0.52\linewidth}
\begin{exemple}[Ejemplo: \ $3(x-2)=x+4$]
\vspace{-6pt}
\begin{alignat*}{2}
3x-6&=x+4 &\quad& \text{\nota{\pas{1}}}\\
3x-x&=4+6 && \text{\nota{\pas{2}}}\\
2x&=10 && \text{\nota{\pas{3}}}\\
x&=\tfrac{10}{2}=\boxed{5} && \text{\nota{\pas{4}}}
\end{alignat*}
\nota{\pas{5}: $3\cdot3=9$ y $5+4=9$ \be{}}
\end{exemple}
\end{minipage}

\vfill
\begin{compte}[El signo menos delante de un paréntesis]
Cambia el signo de \textbf{todos} los términos de dentro: \ $10-(x+3)=10-x-3$ \qquad $-2(x-4)=-2x+8$ \qquad \mal{-2(x-4)=-2x-8}
\end{compte}

\vfill
\begin{exemple}[Ejemplo: \ $4(x+1)-2(x-3)=3x+5$ \ (dos paréntesis y un signo menos)]
$4x+4-2x+6=3x+5$ \ $\to$ \ $4x-2x-3x=5-4-6$ \ $\to$ \ $-x=-5$ \ $\to$ \ $\boxed{x=5}$
\end{exemple}

\vfill
\begin{minipage}[t]{0.48\linewidth}
\begin{exemple}[Ejemplo: \ $5(2x-1)=3(x+4)$]
$10x-5=3x+12$\\[3pt]
$7x=17$ \ $\to$ \ $\boxed{x=\tfrac{17}{7}}$ \ \nota{puede ser una fracción}
\end{exemple}
\end{minipage}\hfill
\begin{minipage}[t]{0.49\linewidth}
\begin{exemple}[Ejemplo: \ $12-(2x+5)=x+1$]
$12-2x-5=x+1$\\[3pt]
$-3x=-6$ \ $\to$ \ $\boxed{x=2}$
\end{exemple}
\end{minipage}

\vfill
\begin{exemple}[Ejemplos]
\begin{minipage}[t]{0.48\linewidth}
$2(x+3)=14$: \ $2x+6=14$, \ \boxed{x=4}\\[4pt]
$6-3(x-1)=0$: \ $6-3x+3=0$, \ \boxed{x=3}
\end{minipage}\hfill
\begin{minipage}[t]{0.49\linewidth}
$3(x+2)-x=2(x+3)$: \ $2x+6=2x+6$\\ \hspace*{1em}\nota{infinitas soluciones: una identidad}\\[4pt]
$2(x-1)=2x+5$: \ $0x=7$ \ \nota{ninguna solución}
\end{minipage}
\end{exemple}

\vfill
% ============================================================ 6.8
\clearpage
\seccio{sec8}{Ecuaciones con denominadores}{Primero nos libramos de las fracciones; después, como siempre.}

\vfill
\begin{idea}
{\centering\Large Multiplicamos \textbf{todos} los términos por el \fmkey{m.c.m.} de los denominadores.\par}
\smallskip Obtenemos una ecuación \textbf{equivalente} sin fracciones y seguimos como en el apartado anterior.
\end{idea}

\vfill
\begin{center}
\begin{tikzpicture}[t/.style={font=\fontsize{24}{24}\selectfont,inner sep=3pt},
  r/.style={font=\fontsize{20}{20}\selectfont,inner sep=3pt},
  l/.style={font=\comic\small,text=fmmuted,anchor=west}]
  \node[t] (a) at (0,0) {$\dfrac{\cA{x+1}}{2}$};
  \node[t] (m) at (1.6,0) {$-$};
  \node[t] (b) at (3.2,0) {$\dfrac{\cB{x-2}}{3}$};
  \node[t] (e) at (4.8,0) {$=$};
  \node[t] (c) at (6.0,0) {$\cK{2}$};
  \node[font=\Large\heavy,fill=fmyellow!60,rounded corners=2mm,inner sep=5pt,anchor=west] (mcm) at (8.0,0) {m.c.m.$(2,3)=6$};
  \node[r] (a2) at (0,-2.5) {$3\,(\cA{x+1})$};
  \node[r] (m2) at (1.6,-2.5) {$-$};
  \node[r] (b2) at (3.2,-2.5) {$2\,(\cB{x-2})$};
  \node[r] (e2) at (4.8,-2.5) {$=$};
  \node[r] (c2) at (6.0,-2.5) {$12$};
  \foreach \u/\v/\t in {a/a2/{6:2=3},b/b2/{6:3=2},c/c2/{2\cdot6=12}} {
    \draw[-{Latex[length=2.5mm]},fmmuted,line width=1.1pt] ([yshift=-2pt]\u.south) -- ([yshift=2pt]\v.north);
    \node[l] at ($(\u.south)!0.5!(\v.north)+(3pt,0)$) {$\t$};}
  \node[font=\comic,text=fmmuted,align=left,anchor=north west] at (8.0,-0.75) {Dividimos el m.c.m. entre cada\\denominador y el resultado\\\textbf{multiplica} al numerador,\\que va \textbf{entre paréntesis}.};
  \node[font=\Large,anchor=west] at (-0.9,-3.9) {$\Rightarrow\ 3x+3-2x+4=12\ \Rightarrow\ x+7=12\ \Rightarrow\ \boxed{x=5}$};
\end{tikzpicture}
\end{center}

\vfill
\begin{minipage}[t]{0.52\linewidth}
\begin{compte}[El menos delante de una fracción]
Afecta a \textbf{todo} el numerador. Por eso lo ponemos siempre \textbf{entre paréntesis}:\\[4pt]
\be{-\dfrac{x-2}{3}\ \to\ -2(x-2)=-2x\cR{+4}}\\[4pt]
\mal{-\dfrac{x-2}{3}\ \to\ -2x-4}
\end{compte}
\end{minipage}\hfill
\begin{minipage}[t]{0.45\linewidth}
\begin{compte}[¡No te olvides!]
Los términos \textbf{sin denominador} también se multiplican por el m.c.m.:\\[4pt]
\be{\ldots=\cK{2}\ \to\ \ldots=12}\qquad \mal{\ldots=2}
\end{compte}
\end{minipage}

\vfill
\begin{minipage}[t]{0.58\linewidth}
\begin{exemple}[Ejemplo: \ $\dfrac{x-1}{4}-\dfrac{x-3}{6}=1$]
\vspace{-6pt}
\begin{alignat*}{2}
&3(x-1)-2(x-3)=12 &\qquad& \text{\nota{m.c.m.$(4,6)=12$}}\\
&3x-3-2x+6=12\\
&x+3=12\ \to\ \boxed{x=9}
\end{alignat*}
\end{exemple}
\end{minipage}\hfill
\begin{minipage}[t]{0.39\linewidth}
\begin{truc}[¡No pasa nada!]
La solución puede ser una \textbf{fracción}:\\[4pt]
$-2x=-13\ \to\ x=\dfrac{13}{2}$\\[3pt]
también es una solución válida.
\end{truc}
\end{minipage}

\vfill
\begin{exemple}[Ejemplo: \ $\dfrac{x}{3}+\dfrac{x}{2}=10$]
\begin{minipage}[c]{0.55\linewidth}
m.c.m.$(3,2)=6$: \quad $2x+3x=60\ \to\ 5x=60\ \to\ \boxed{x=12}$
\end{minipage}\hfill
\begin{minipage}[c]{0.42\linewidth}
\nota{Comprobación: $\dfrac{12}{3}+\dfrac{12}{2}=4+6=10$ \be{}}
\end{minipage}
\end{exemple}

\vfill
% ============================================================ 6.9
\clearpage
\seccio{sec9}{Resolución de problemas con ecuaciones}{Traducimos el enunciado a lenguaje algebraico.}

\passosproblema{\textbf{Elegimos}\\la incógnita}{\textbf{Planteamos}\\la ecuación}

\vfill
\begin{minipage}[c]{0.58\linewidth}\centering\small
\taularod{>{\hspace{2pt}}l<{\hspace{2pt}}>{\hspace{2pt}}c<{\hspace{2pt}}}{
\rowcolor{fmblue}\color{white}\heavy Enunciado & \color{white}\heavy Lenguaje algebraico\\
Tres números consecutivos & $x,\ x+1,\ x+2$\\
\rowcolor{fmsurface}Dos pares consecutivos & $2x,\ 2x+2$\\
Edad dentro de 5 años / hace 3 & $x+5\ /\ x-3$\\
\rowcolor{fmsurface}Dos tercios de un número & $\frac23x$\\
Precio con un 20\,\% de descuento & $0{,}8\,x$\\
\rowcolor{fmsurface}Dos cantidades que suman 50 & $x,\ 50-x$
}
\end{minipage}\hfill
\begin{minipage}[c]{0.39\linewidth}
\begin{compte}[Comprueba e interpreta]
La solución debe tener sentido: una edad no puede ser negativa y un número de personas debe ser natural. Responde con una frase.
\end{compte}
\end{minipage}

\vfill
\begin{minipage}[t]{0.48\linewidth}
\begin{exemple}[Edades: \ padre de 42 e hija de 12. ¿Cuándo tendrá el triple?]
Dentro de $x$ años: \ $42+x=3(12+x)$\\[3pt]
$42+x=36+3x\ \to\ 6=2x\ \to\ x=3$\\[3pt]
\boxed{\text{dentro de 3 años}} \ \nota{45 y 15}
\end{exemple}
\end{minipage}\hfill
\begin{minipage}[t]{0.49\linewidth}
\begin{exemple}[Fracciones: \ gasta $\frac25$ y $\frac14$; le quedan 105 €]
$x-\dfrac25x-\dfrac14x=105$ \ \nota{por 20}\\[3pt]
$20x-8x-5x=2100\ \to\ 7x=2100$\\[3pt]
\boxed{x=300\text{ €}}
\end{exemple}
\end{minipage}

\vfill
\begin{minipage}[t]{0.48\linewidth}
\begin{exemple}[Mezclas: \ 50 kg a 9 €/kg con café de 12 €/kg y de 7 €/kg]
$x$ kg del de 12 €; \ $50-x$ del de 7 €:\\[3pt]
$12x+7(50-x)=9\cdot50\ \to\ 5x=100$\\[3pt]
\boxed{20\text{ kg y }30\text{ kg}}
\end{exemple}
\end{minipage}\hfill
\begin{minipage}[t]{0.49\linewidth}
\begin{exemple}[Móviles: \ 120 km, uno hacia el otro a 25 y a 35 km/h]
En $x$ horas hacen $25x$ y $35x$ km:\\[3pt]
$25x+35x=120\ \to\ 60x=120$\\[3pt]
\boxed{x=2\text{ horas}}
\end{exemple}
\end{minipage}

\vfill
\begin{exemple}[Ejemplos]
Tres pares consecutivos que suman 114: \ $2x+(2x+2)+(2x+4)=114\ \to\ x=18$: \ \boxed{36,\ 38,\ 40}\\[4pt]
Depósito lleno a $\frac34$; si sacamos 120 L, queda a $\frac13$: \ $\frac34x-120=\frac13x\ \to\ 9x-1440=4x\ \to\ \boxed{x=288\text{ L}}$
\end{exemple}

\vfill

\vfill
% ============================================================ resum
\resumtitol
\begin{tcbraster}[raster columns=2,raster equal height=rows,raster column skip=14pt,raster row skip=24pt,raster before skip=14pt]
\rtile{6.1}{Lenguaje algebraico}{Letras para números desconocidos.\\[1ex] El doble, $2x$; el siguiente, $x+1$.}
\rtile{6.2}{Valor numérico}{Sustituimos y operamos.\\[1ex] Negativos, entre paréntesis.}
\rtile{6.3}{Monomios}{Coeficiente, parte literal y grado.}
\rtile{6.4}{Operaciones}{Sumamos solo los semejantes.\\[1ex] Multiplicar: coeficientes y letras.}
\rtile{6.5}{Ecuaciones}{Igualdad con incógnita.\\[1ex] Solución: el valor que la hace cierta.}
\rtile{6.6}{Resolución}{Transponemos, reducimos, despejamos y comprobamos.}
\rtile{6.7}{Paréntesis}{Primero los quitamos.\\[1ex] Un $-$ delante cambia todos los signos.}
\rtile{6.8}{Denominadores}{Multiplicamos por el mcm.}
\rtile{6.9}{Problemas}{Elegimos $x$ y planteamos.\\[1ex] Resolvemos y comprobamos.}
\end{tcbraster}
\end{document}
